\documentclass[11pt,a4paper]{article}

\usepackage[margin=2cm]{geometry}
\usepackage{titlesec}
\usepackage{enumitem}
\usepackage{hyperref}
\usepackage{xcolor}

% Colors
\definecolor{linkcolor}{RGB}{0,70,130}

% Hyperref setup
\hypersetup{
    colorlinks=true,
    linkcolor=linkcolor,
    urlcolor=linkcolor
}

% Section formatting
\titleformat{\section}{\large\bfseries\uppercase}{}{0em}{}[\titlerule]
\titlespacing{\section}{0pt}{12pt}{6pt}

% Remove page numbers
\pagenumbering{gobble}

% Custom commands
\newcommand{\entry}[4]{
    \textbf{#1} \hfill #2 \\
    \textit{#3} \hfill \textit{#4}
}

\begin{document}

% Header
\begin{center}
    {\LARGE\bfseries Vinícius Bregoli} \\[4pt]
    Desenvolvedor de Software \\[6pt]
    \href{https://www.viniciusbregoli.dev}{viniciusbregoli.dev} \enspace$\cdot$\enspace
    \href{mailto:me@viniciusbregoli.dev}{me@viniciusbregoli.dev} \enspace$\cdot$\enspace
    Curitiba, Paraná \\[2pt]
    \href{https://linkedin.com/in/viniciusbregoli/}{linkedin.com/in/viniciusbregoli} \enspace$\cdot$\enspace
    \href{https://github.com/viniciusbregoli}{github.com/viniciusbregoli}
\end{center}

\section{Habilidades Técnicas}
\noindent\textbf{Linguagens:} Python, Rust, JavaScript, C/C++, SQL, HTML, CSS \\
\textbf{Frameworks e Bibliotecas:} Node.js, Pytorch, OpenCV \\
\textbf{Ferramentas e Plataformas:} Git, Linux, Virtualização \\
\textbf{Áreas:} Desenvolvimento Full-Stack, Machine Learning, Deploy, Infraestrutura de TI

\section{Experiência Profissional}

\entry{Desenvolvedor ML/IA}{Kinebot}{Curitiba, Brasil}{Set 2025 -- Presente}
\begin{itemize}[nosep, leftmargin=*]
    \item Desenvolvimento de soluções de visão computacional para estimativa de pose humana
    \item Implementação e ajuste fino de aplicações baseadas em LLMs
    \item Construção de pipelines de ML principalmente em Python usando PyTorch e OpenCV
\end{itemize}

\vspace{8pt}

\entry{Estagiário -- Desenvolvimento de Software}{Exati Tecnologias}{Curitiba, Brasil}{2022 -- 2023}
\begin{itemize}[nosep, leftmargin=*]
    \item Desenvolvimento e manutenção de aplicações web utilizando Vue.js
    \item Integração de aplicações frontend com serviços backend usando Java e Node.js
    \item Criação e gerenciamento de bancos de dados PostgreSQL
    \item Colaboração utilizando Git/GitLab para controle de versão e YouTrack para gestão de tarefas
\end{itemize}

\vspace{8pt}

\entry{Estagiário -- Infraestrutura de TI}{Siemens Energy Brasil}{Curitiba, Brasil}{2021 -- 2022}
\begin{itemize}[nosep, leftmargin=*]
    \item Gerenciamento de infraestrutura de TI e ambientes de virtualização
    \item Suporte técnico remoto em ambiente corporativo de alta demanda
    \item Gestão de chamados e tickets
\end{itemize}

\section{Formação Acadêmica}

\entry{Ciência da Computação e IA (Intercâmbio)}{Technische Hochschule Ingolstadt}{Ingolstadt, Alemanha}{Out 2024 -- Ago 2025}

\vspace{8pt}

\entry{Bacharelado em Engenharia de Computação}{Pontifícia Universidade Católica do Paraná (PUCPR)}{Curitiba, Brasil}{2020 -- 2025}

\section{Projetos e Pesquisa}

\textbf{Pesquisa em Digital Twins -- Parceria Airbus} \hfill 2021 \\
Programa de Bolsas de Iniciação em Desenvolvimento Tecnológico e Inovação, com foco em digital twins e suas aplicações em projetos aeroespaciais.

\vspace{8pt}

\textbf{Equipe de Robótica Móvel PUCPR} \hfill 2020 -- 2023 \\
Contribuição no desenvolvimento de robôs seguidores de linha, sumô, trekking e combate. Participação em competições como RoboCore Experience (RCX) e IRONCup.

\section{Idiomas}

\noindent Inglês (Fluente) \enspace$\cdot$\enspace Português (Nativo) \enspace$\cdot$\enspace Espanhol (Intermediário) \enspace$\cdot$\enspace Alemão (Básico)

\end{document}
